\section*{Chapitre \ref{ch:rc:regulatory-capital-for-trading-books} --- Capital réglementaire des portefeuilles de négociation}

\begin{solution}{exo:rc:regulatory-capital-for-trading-books:1}
$\max(e^{-0.03\times8/2},0.40) = e^{-0.12} = 88.7~\%$~; haute $\min(1.25\times0.887,1) = 100~\%$~; basse
$\max(2\times0.887-1,0.75\times0.887) = 77.4~\%$.
\end{solution}

\begin{solution}{exo:rc:regulatory-capital-for-trading-books:2}
Le portefeuille est long de la maturité à dix ans et court de celle à deux ans~: les positions se compensent grâce à leur corrélation positive,
et une corrélation plus faible signifie moins de compensation. Il en va de même pour les compartiments de change long en euro et
court en yen. Les charges de vega et de courbure restant inchangées, c'est le scénario de corrélation basse qui ajoute le plus.
\end{solution}

\begin{solution}{exo:rc:regulatory-capital-for-trading-books:3}
Zone orange~: la corrélation passe (au-dessus de 0.80), mais la distance de distribution (0.10) se situe entre 0.09 et
0.12. Le desk conserve le modèle interne mais acquitte la surcharge.
\end{solution}

\begin{solution}{exo:rc:regulatory-capital-for-trading-books:4}
L'approche standardisée impute la courbure du straddle (un choc de $\pm10.6~\%$ sur EURUSD) et son vega à un poids de
100~\%, des paramètres calibrés pour n'importe quelle banque~; le modèle interne mesure l'\omterm{def:rc:market-risk-measures:es}{expected shortfall} à dix jours du portefeuille réel,
où les mouvements d'EURUSD restent bien plus faibles, même l'année de tensions, et la diversification
s'applique.
\end{solution}

\begin{solution}{exo:rc:regulatory-capital-for-trading-books:5}
Non~: il compte plus de 24 observations, mais les huit derniers mois comportent des périodes de 90 jours sans aucune, si bien que
le critère (1) échoue, et 30 est inférieur aux 100 du critère (2). Il est non modélisable.
\end{solution}

\begin{solution}{exo:rc:regulatory-capital-for-trading-books:6}
L'\omterm{def:rc:market-risk-measures:es}{expected shortfall} diversifié repose sur les corrélations entre classes calculées par le modèle, qui peuvent se rompre en
période de tensions, alors que la somme sur les classes suppose qu'elles sont nulles. Pondérer les deux à parts égales limite
le bénéfice en capital de la diversification entre classes.
\end{solution}

\begin{solution}{exo:rc:regulatory-capital-for-trading-books:7}
$\sqrt{5^2+\bigl(3\sqrt{(20-10)/10}\bigr)^2} = \sqrt{34} = 5.83$.
\end{solution}

\begin{solution}{exo:rc:regulatory-capital-for-trading-books:8}
Un chiffre plus bas n'est pas meilleur~: le modèle doit réussir le backtesting et le \omterm{def:rc:regulatory-capital-for-trading-books:pla}{test d'attribution du P\&L}, desk par desk, sa calibration
en période de tensions et ses facteurs non modélisables ajoutent du capital, et l'approche standardisée demeure un indice de référence
(et, dans certaines juridictions, un plancher). Un montant faible assorti d'un test d'attribution échoué est
un élément à charge contre le modèle.
\end{solution}

\begin{solution}{pb:rc:regulatory-capital-for-trading-books:1}
\textbf{1.} 68.29~millions de USD, dans le scénario de corrélation basse.
\textbf{2.} 25.56~millions de USD~: 1.5 fois un IMCC de 17.04~millions de USD, à partir d'un ES en période de tensions de 14.41~millions de USD
pour le portefeuille, et de 5.67 et 14.01~millions de USD pour les taux et le change.
\textbf{3.} Courbure de change, 39.23~millions de USD, provenant du straddle vendu.
\textbf{4.} Il contient la hausse du dollar et la baisse de l'euro sous la parité, les plus forts mouvements d'EURUSD
de l'échantillon, auxquels le straddle vendu et la position en euros sont le plus sensibles.
\textbf{5.} Rien à signaler dans le modèle interne ici, puisque la volatilité n'est pas un facteur modélisé~; dans un modèle qui la
prend en compte, un horizon de 40 jours augmenterait l'ES porté par le vega.
\textbf{6.} 0.703 et 0.044.
\textbf{7.} Zone orange~: la corrélation est comprise entre 0.70 et 0.80~; la distance passe.
\textbf{8.} Une surcharge de $0.5\times(68.29-25.56) = \text{USD}~21.36$ millions~; capital de 46.93~millions de USD.
\textbf{9.} L'exigence standardisée, 68.29~millions de USD.
\textbf{10.} Les modèles delta--gamma et sans yen passent (vert)~; le modèle sans deux ans est orange
en raison de sa distance (0.100).
\textbf{11.} Le proxy fait bouger les deux maturités ensemble, si bien qu'il manque les jours où la courbe pivote~: l'ordre des
P\&L quotidiens (rangs) est brouillé, tandis que la dispersion de la distribution reste proche.
\textbf{12.} Des observations de prix réels de swaps ou d'obligations à dix ans~: au moins 24 par an, sans lacune de 90 jours
comportant moins de quatre.
\textbf{13.} En partie~: une régression recale l'amplitude du mouvement mais ne voit toujours pas le pivot de la courbe~; seul un facteur distinct
restaure les rangs.
\textbf{14.} Il est temps de collecter des données et de corriger le modèle avant que le test n'affecte le capital.
\textbf{15.} Comparer le coût de la correction à la valeur actuelle de la surcharge sur l'horizon du desk, et le
risque de passer au rouge.
\textbf{16.} Le backtesting contrôle le nombre de dépassements en queue~; l'attribution vérifie que le modèle de risque voit chaque jour les mêmes facteurs
explicatifs du P\&L. Un modèle peut réussir l'un et échouer l'autre.
\textbf{17.} Pour améliorer les modèles de risque, ou pour manipuler les statistiques en choisissant des proxys qui préservent les rangs.
\textbf{18.} Lorsque les règles et les multiplicateurs diffèrent, une banque peut enregistrer ses opérations dans la juridiction au
traitement le plus léger, ou faire tourner des modèles différents.
\textbf{19.} Cas-type~: \emph{le desk qui a échoué au test d'attribution}~: capital du modèle interne 25.56~millions de USD,
standardisé 68.29~millions de USD~; Spearman 0.703 et KS 0.044 le placent en zone orange, et le capital passe à 46.93
millions de USD.
\textbf{20.} Que le modèle de risque et les systèmes de valorisation voient le même monde.
\end{solution}

\begin{solution}{iq:rc:regulatory-capital-for-trading-books:1}
La VaR et la VaR en période de tensions, à 99~\%, sont remplacées par l'\omterm{def:rc:market-risk-measures:es}{expected shortfall} à 97.5~\%, avec des \omterm{def:rc:regulatory-capital-for-trading-books:ima}{horizons de liquidité} et une
calibration en période de tensions~; une \omterm{def:rc:regulatory-capital-for-trading-books:sbm}{méthode fondée sur les sensibilités}~; un agrément au niveau du desk avec backtesting et
\omterm{def:rc:regulatory-capital-for-trading-books:pla}{test d'attribution du P\&L}~; des \omterm{def:rc:regulatory-capital-for-trading-books:nmrf}{facteurs de risque non modélisables}~; une délimitation plus stricte du \omterm{def:rc:regulatory-capital-for-trading-books:books}{portefeuille de négociation}.
\iqlookfor{ES, \omterm{def:rc:regulatory-capital-for-trading-books:ima}{horizons de liquidité}, SBM et tests au niveau du desk.}
\end{solution}

\begin{solution}{iq:rc:regulatory-capital-for-trading-books:2}
PV01 par devise et par maturité~; pondérations de risque par maturité (divisées par $\sqrt2$ pour les devises majeures)~;
agrégation au sein de chaque devise avec les corrélations entre maturités, puis entre devises à 50~\%~; les trois
scénarios de corrélation~; la courbure, à partir de chocs parallèles, s'il y a des options.
\iqlookfor{les étapes et les paramètres.}
\end{solution}

\begin{solution}{iq:rc:regulatory-capital-for-trading-books:3}
Il compare le P\&L du modèle de risque à celui du front office, par corrélation de rang et par une
distance entre distributions. Les desks échouent à cause de \omterm{def:rc:market-risk-measures:factor}{facteurs de risque} manquants ou traités par proxy, de sources de données ou d'heures de capture différentes,
et de réévaluations simplifiées, et des risques de base.
\iqlookfor{les indicateurs et les causes courantes.}
\end{solution}

\begin{solution}{iq:rc:regulatory-capital-for-trading-books:4}
L'ES capte l'ampleur des pertes en queue et il est cohérent~; les \omterm{def:rc:regulatory-capital-for-trading-books:ima}{horizons de liquidité} reflètent le fait que certaines positions prennent
des semaines pour se dénouer, ce qu'une VaR uniforme à dix jours ignore~; la calibration en période de tensions évite que le capital baisse
sur les marchés calmes.
\iqlookfor{queue, liquidité et procyclicité.}
\end{solution}

\begin{solution}{iq:rc:regulatory-capital-for-trading-books:5}
Par l'économie de capital face au coût des modèles, des données et des contrôles, et par la probabilité de réussir
les tests de backtesting, d'attribution et d'éligibilité~; en premier lieu, les desks aux risques liquides et bien modélisés.
\iqlookfor{coût-bénéfice et risque lié aux tests.}
\end{solution}

\begin{solution}{iq:rc:regulatory-capital-for-trading-books:6}
De longues séries de prix réels par \omterm{def:rc:market-risk-measures:factor}{facteur de risque}, avec le décompte des observations~; des données de marché cohérentes pour
le front office et la fonction risque~; le \omterm{def:rc:regulatory-capital-for-trading-books:hpl}{P\&L hypothétique} et le \omterm{def:rc:regulatory-capital-for-trading-books:hpl}{P\&L théorique du risque} par desk~; une infrastructure de réévaluation complète pour l'ES sur de nombreux
scénarios et sous-ensembles de liquidité~; recherche de périodes de tensions~; reporting et pistes d'audit.
\iqlookfor{données, réévaluation et contrôles.}
\end{solution}
